\section{Epistemic conditions for Nash equilibrium}
\label{sec:annex-p17-epistemic-nash-conditions}

\begin{annexitem}{annex:p17.ab.measurability}{thm:p17-ab-measurability}{The source's measurable-space extension made explicit}
\begin{proof}
Fix \(s_i\). The insertion map from \(S^{-i}\) to \(S\), which
inserts \(s_i\) as coordinate \(i\), is measurable: the inverse
image of a measurable rectangle is a rectangle or empty.
The sets whose inverse images are measurable form a sigma-field,
so the assertion extends from rectangles to the product sigma-field.
Pushing \(b_i(s_i)\) forward by this map gives a countably additive
probability \(p_i(\cdot;s_i)\). Measurability in \(s_i\) is the
explicit kernel hypothesis, not an assertion that disintegrations
exist on arbitrary measurable spaces.

Action-profile equality events are finite intersections of measurable
action fibers. Thus each coordinate of \(\phi^i\) is measurable.
Each payoff array has finitely many measurable real coordinates;
equalities with a fixed array are finite intersections of their
level sets. Each expected-payoff difference defining \(R_i\) is
a finite sum of products of measurable, finite real functions.
Choosing the actual own action can be written as a finite sum over
its measurable action fibers. Hence \(R_i\) is measurable and depends
only on \(s_i\). No bound uniform over all types is needed: the
proof never integrates these payoffs against the common prior.

The kernel assumption makes \(K_iE\) measurable for every measurable
\(E\). Finite intersections and induction give every \(K^rE\),
and a countable intersection gives \(CKE\).
Taking \(B=S_i\) in (CP) gives the averaging identity. Integrating
the stated equality of probabilities then gives the proportional
identity, including its almost-everywhere form. Only bounded
probability functions are integrated.
\end{proof}
\end{annexitem}

\begin{annexitem}{annex:p17.ab.introspection}{thm:p17-ab-introspection}{Own-type events, countable operations and finite computation}
\begin{proof}
If \(D\) depends only on coordinate \(i\), insertion of the fixed
\(s_i\) puts every possible opponent tuple either inside \(D\)
or outside it. Therefore \(p_i(D;s_i)=1_D(s)\), proving
\(K_iD=D\). The probability-level event in the statement is of
this form because its probability function is measurable and
depends only on \(s_i\).

For countably many events \(E_k\), if each has probability one,
their intersection has probability one: its complement is a
countable union of null events, whose probability is at most the
sum of their zero probabilities. Conversely a probability-one
intersection is contained in each event, so each event has
probability one. This proves preservation of countable
intersections. The same containment argument proves monotonicity.

Write \(K^rE\) for the successive iterations. Since
\(K^{r+1}E\subseteq K_i(K^rE)\),
\[
 CKE\subseteq\bigcap_{r\geq1}K_i(K^rE)
       =K_i\!\left(\bigcap_{r\geq1}K^rE\right)=K_iCKE.
\]
For the conjecture event \(D=[\phi=\varphi]\), knowledge of \(D\)
implies knowledge of the own-coordinate event
\([\phi^i=\varphi^i]\). Own-type certainty therefore implies
\(\phi^i(s)=\varphi^i\) whenever \(s\in K^1D\).
Applying this for all \(i\) gives \(K^1D\subseteq D\).
Since \(CKD\subseteq K^1D\), its assertion follows too.

For the finite computation, draw an edge \(s\to t\) when some
player's extended probability at \(s\) gives \(t\) positive mass.
Membership in \(K^1E\) is exactly membership of every successor in
\(E\). Induction makes \(K^rE\) the condition that every endpoint
of every length-\(r\) path lies in \(E\).
Thus \(CKE\) means that every state reachable by a path of positive
length lies in \(E\). The starting state is not automatically
included unless reachable by such a path.
Every shortest positive-length path between distinct states has
length at most \(|S|-1\); a shortest return path has length at most
\(|S|\). Consequently it suffices to check lengths
\(1,\ldots,|S|\). This proves the finite procedure even when
successive \(K^rE\) themselves are not decreasing.
\end{proof}
\end{annexitem}

\begin{annexitem}{annex:p17.ab.preliminary}{thm:p17-ab-preliminary}{Known actions and optimality}
\begin{proof}
Knowledge of \([a=\alpha]\) implies, for player \(i\), knowledge
of its own action equality. Own-type certainty gives
\(a_i(s_i)=\alpha_i\); hence the actual profile is \(\alpha\).
The same probability-one event puts the entire conjecture on
\(\alpha^{-i}\). Rationality therefore says
\[
 g_i(s_i)(\alpha_i,\alpha^{-i})
       \geq g_i(s_i)(d_i,\alpha^{-i})\quad(d_i\in A_i).
\]
These are the pure Nash inequalities for the actual payoff array
of each player, proving the result. No belief about the other
players' rationality or payoff arrays was used.
\end{proof}
\end{annexitem}

\begin{annexitem}{annex:p17.ab.supported.actions}{thm:p17-ab-supported-actions}{A positive event intersected with three certain events}
\begin{proof}
Choose a state \(s\) in the stated intersection and a player
\(i\ne j\) whose specified conjecture gives positive marginal
probability to \(c_j\). Introspection gives
\(\phi^i(s)=\varphi^i\), so \(p_i([a_j=c_j];s_i)>0\).
The events \(R\), \([g=g^*]\) and \([\phi=\varphi]\) each have
probability one under this same measure. Removing their
complements from \([a_j=c_j]\) removes probability zero.
Their fourfold intersection therefore has positive probability
and is nonempty.

At a state in that intersection, player \(j\) chooses \(c_j\),
has payoff array \(g_j^*\), has conjecture \(\varphi^j\),
and is rational. Its defining rationality inequalities are exactly
the asserted best-response inequalities. This argument only
requires that a positive-measure event be nonempty. It never
selects a positive-probability individual state, so it remains
valid for nonatomic type spaces.
\end{proof}
\end{annexitem}

\begin{annexitem}{annex:p17.ab.two.player}{thm:p17-ab-two-player}{Theorem A on the full type domain}
\begin{proof}
Set \(\sigma_1=\varphi^2\) and \(\sigma_2=\varphi^1\).
Every action with positive \(\sigma_1\) probability is a best
response to \(\sigma_2\), by the preceding lemma; reversing the
players gives the corresponding assertion for \(\sigma_2\).
A mixture supported on payoff-maximizing actions has the same
payoff as their common maximum, because its nonnegative weights
sum to one. Every other mixture averages payoffs no larger
than that maximum. Thus neither player can improve by a mixed
deviation. These are the equilibrium conditions of
Definition~\ref{def:p17-nash-equilibrium}, equivalently
Theorem~\ref{thm:p17-best-response-geometry}.
For the distinction from actual pure play, take the two-player
matching-pennies payoffs given by the first two coordinates of
Definition~\ref{def:p17-ab-chain-game}. Give each player one type
per action, with that action, the fixed payoff array, and fair
belief on opponent types. Their independent uniform prior is
common. All types are rational, and the game and conjectures
are constant, so all Theorem A hypotheses hold everywhere.
At the positive-prior state \((H,h)\), column can improve from
payoff 0 to 1 by choosing \(t\). Thus the actual profile is
not a pure equilibrium although both conjectures are the
equilibrium fair mixtures.
\end{proof}
\end{annexitem}

\begin{annexitem}{annex:p17.ab.type.bridge}{thm:p17-ab-type-bridge}{Both finite encodings and their exact hypotheses}
\begin{proof}
The stated pushforward is a probability on the product of opponent
choice/type pairs. Its action marginal is the original conjecture,
because both marginalizations sum the same original type weights.
The fixed payoff functions then give identical expected payoffs
for every own action in both models. A source action is rational
exactly when it is optimal at the corresponding preceding-model type.

In the two-player case, source mutual rationality puts probability
one on rational opponent choice/type pairs. Source mutual knowledge
of the full conjecture profile makes the conjecture of every
positively conjectured opponent type equal to the specified one.
Introspection makes that specified conjecture the actual one.
These are precisely the two hypotheses of the preceding theorem.

For the reverse construction, type \((c_i,t_i)\) has action \(c_i\),
fixed own payoff function, and the opponent-pair law \(b_i(t_i)\)
already supplied by the preceding model. Thus its action conjecture
is \(q_i(t_i)\), independently of its own first coordinate.
At each actual \(t_i\), the finite nonempty action set has an
expected-payoff maximizer; choose one as \(c_i\).
The resulting source actual own action is rational.
Belief in the opponent's rationality in the preceding model makes
all positively conjectured opponent pairs rational in the source
model. Since the own coordinate remains fixed, each player assigns
probability one to all-player rationality. The specified correctness
of opponent conjectures, together with fixed own conjecture, gives
mutual knowledge of the full actual conjecture profile.
The game is constant everywhere. Hence Theorem A applies at the
constructed state. No common prior was inserted by this encoding.
\end{proof}
\end{annexitem}

\begin{annexitem}{annex:p17.ab.illustration}{thm:p17-ab-illustration}{All types in Figures 1 and 2}
\begin{proof}
Every row and column type has marginal probability \(1/3\),
so conditioning splits equally between its two occupied cells.
Row types \(C_1,D_1\) each conjecture \(c,d\) with probabilities
\(1/2,1/2\), yielding expected payoffs \((1,1/2)\) to \((C,D)\).
Thus \(C_1\) is rational and \(D_1\) is not. Type \(D_2\)
conjectures \(d\) with certainty, yielding \((0,1)\), and is rational.
Column types \(c_1,d_1\) analogously yield payoffs \((1,1/2)\)
to \((c,d)\), so \(c_1\) is rational and \(d_1\) is not.
Type \(d_2\) conjectures \(D\) with certainty, yields \((0,1)\),
and is rational.

At \((D_2,d_2)\), row's two possible column types both choose
\(d\), and column's two possible row types both choose \(D\).
Together with certainty of their own choices, this proves mutual
knowledge of the actual action profile. Row assigns probability
\(1/2\) to irrational column type \(d_1\); column assigns
\(1/2\) to irrational row type \(D_1\). Mutual rationality fails.

At the state \((D_2,d_1)\), which row considers possible,
column assigns positive probability to row type \(C_1\), whose
action is \(C\). Thus column does not know the action profile
\((D,d)\) there. Row therefore does not know that the action
profile is mutually known at the original state. Second-order,
and hence common, knowledge fails. The direct payoff comparison
at \((D,d)\) gives \(1>0\) for both players, verifying its
pure Nash property independently.
\end{proof}
\end{annexitem}

\begin{annexitem}{annex:p17.ab.local.average}{thm:p17-ab-local-average}{Averaging without an unjustified measurable projection}
\begin{proof}
Fix a type \(t_i\). If there is a state in \(H\) with that
coordinate, the inclusion \(H\subseteq K_iH\) gives
\(p_i(H;t_i)=1\), and the hypothesis on \(E\) gives
\(p_i(E;t_i)=\lambda\). Consequently
\(p_i(E\cap H;t_i)=\lambda=\lambda p_i(H;t_i)\).
If there is no such state, the section of \(H\) at \(t_i\)
is empty, so both \(p_i(H;t_i)\) and \(p_i(E\cap H;t_i)\)
are zero. The same equality therefore holds at every type.
Integrating it using the common-prior averaging identity proves
\(P(E\cap H)=\lambda P(H)\).

No projected set is integrated or assumed measurable in this
argument. In fact the set of types with nonempty \(H\)-section
equals the measurable set \(\{t_i:p_i(H;t_i)=1\}\) under the
absorption hypothesis. Thus the source's projection reasoning
also obtains the measurability it needs here, rather than
assuming it for general measurable projections.
\end{proof}
\end{annexitem}

\begin{annexitem}{annex:p17.ab.product}{thm:p17-ab-product}{The complete finite factorization argument}
\begin{proof}
For one coordinate the claim is the definition of its marginal.
For two it is precisely the premise. Suppose it holds for
\(n-1\) coordinates. Sum the assumed equality for coordinate
\(i<n\) over coordinate \(n\). This shows that in the marginal
law on the first \(n-1\) coordinates, coordinate \(i\) is
independent of all the other coordinates in that marginal.
The induction hypothesis factors that marginal as
\(\prod_{i<n}Q_i\). The given independence of coordinate
\(n\) from the rest now gives
\[
 Q(c)=Q_n(c_n)Q_{-n}(c^{-n})=\prod_{i=1}^nQ_i(c_i).
\]
Every operation is summation of finite nonnegative masses;
no division by a marginal, positive or otherwise, is used.
\end{proof}
\end{annexitem}

\begin{annexitem}{annex:p17.ab.many.player}{thm:p17-ab-many-player}{The original common-prior proof with all denominators}
\begin{proof}
Let \(D=[\phi=\varphi]\), \(F=CKD\). Assumption (AB) implies
\(P(F)>0\), so the stated \(Q\) is a probability on the finite
action product. Fix \(i\) and \(c_i\), and put
\(H=[a_i=c_i]\cap F\).
Player \(i\) knows its own action, and \(F\subseteq K_iF\);
intersection preservation gives \(H\subseteq K_iH\).
On \(H\), the inclusion \(F\subseteq D\) makes
\[
 p_i([a^{-i}=c^{-i}];s_i)=\varphi^i(c^{-i}).
\]
The absorbing-event lemma yields
\[
 P([a=c]\cap F)=\varphi^i(c^{-i})P([a_i=c_i]\cap F).
\]
Divide by \(P(F)>0\), the only denominator so far, to obtain
\[
                   Q(c)=\varphi^i(c^{-i})Q_i(c_i).
 \tag{1}
\]
Summing (1) over \(c_i\) gives \(Q_{-i}=\varphi^i\), because
\(\sum_{c_i}Q_i(c_i)=1\). Hence each coordinate is independent
of its complement under \(Q\). The product lemma gives
\(Q=\bigotimes_jQ_j\).

Set \(\sigma_j=Q_j\). For every \(i\ne j\), the \(j\)-marginal
of \(\varphi^i=Q_{-i}\) is \(\sigma_j\). This proves agreement,
and the factorization proves independence of each full conjecture.
Up to this point only (CP) and \(P(F)>0\) have been used;
this proves the separately stated agreement/factorization claim.

The event in (AB) is nonempty. At any state in that event, common
conjecture knowledge implies mutual conjecture knowledge.
The supported-action lemma therefore applies to \(g^*,\varphi\).
If \(\sigma_j(c_j)>0\), choose any \(i\ne j\); its conjecture
has this same positive marginal. Thus \(c_j\) is a best response
to \(\varphi^j=\bigotimes_{k\ne j}\sigma_k\).
The finite Nash support criterion proves equilibrium.

Finally, for every \(j,c_j\) with \(\sigma_j(c_j)=0\),
\(P([a_j=c_j]\cap F)=P(F)Q_j(c_j)=0\).
There are finitely many such actions, so their union is null.
This is the asserted almost-sure support statement. In a finite
model a positive-mass state cannot lie in that union. At an
arbitrary null state the argument gives no such conclusion.
It also shows why merely nonempty \(F\) would not license
the division in (1).
\end{proof}
\end{annexitem}

\begin{annexitem}{annex:p17.ab.marginal.example}{thm:p17-ab-marginal-example}{All conditional laws and deviations in Example 5.1}
\begin{proof}
Each type has prior marginal \(1/2\), and its conditional
distribution splits equally between its two occupied cells.
The opponent laws are
\[
 \begin{array}{c|l}
 U&(LW+RE)/2\\D&(RW+LE)/2\\
 L&(UW+DE)/2\\R&(DW+UE)/2\\
 W&(UL+DR)/2\\E&(UR+DL)/2 .
 \end{array}
\]
Row type \(U\) obtains payoff 1 from \(U\) and \(1/2\) from
\(D\); row type \(D\) obtains 1 from \(D\) and 0 from \(U\).
Each column type obtains 1 from its own lettered action
and 0 from the other action. Each third-player type likewise
obtains 1 from its own action and 0 from the other.
Every type is therefore rational, including at zero-prior
states of the full Cartesian type space.

Every displayed law has fair marginals on each opponent.
Thus game, rationality and the specified individual marginals
hold everywhere, and every knowledge iteration preserves them.
For the product of those fair marginals, row's \(U\) pays 1
at \(LW,RE\) and 0 at the other two equally likely pairs,
giving \(1/2\). Row's \(D\) pays 1 at \(LW,RW,LE\) and
0 at \(RE\), giving \(3/4\). The proposed fair row mixture
has support on the inferior \(U\), so the profile is not Nash.

At \((U,L,W)\), column type \(L\) assigns probability \(1/2\)
to row type \(D\), whose full conjecture differs from that
of row type \(U\). Third type \(W\) does the same.
They therefore do not know row's actual full conjecture.
\end{proof}
\end{annexitem}

\begin{annexitem}{annex:p17.ab.chain.example}{thm:p17-ab-chain-example}{The three boxes, equilibrium and arbitrary finite depth}
\begin{proof}
In the three-box model the total mass is
\(4(9+3+1)/52=1\). Every row or column type has two cells
of equal mass; it therefore assigns a fair distribution to
the opposing matching-pennies action. Both of its actions
give expected payoff \(1/2\), so every such type is rational.

Type \(W_1\) has the two diagonal cells in box 1 with weight
9 each and the two off-diagonal cells in box 2 with weight
3 each. Its probabilities on matching action pairs \(Hh,Tt\)
are \(9/24=3/8\) each, and on mismatching pairs \(Ht,Th\)
are \(3/24=1/8\) each. Type \(W_2\) similarly obtains
\(3/8,3/8,1/8,1/8\) from weights \(3,3,1,1\).
Their payoff from \(W\) is \(3(3/4)=9/4>2\).
Type \(W_3\) sees only the two diagonal pairs and receives
3 from \(W\), versus 2 from \(E\).
Type \(E_1\) sees only off-diagonal pairs and receives
2 from \(E\), versus 0 from \(W\).
Thus every type is rational, and game and rationality hold
throughout the full type product.

At \((H_2,h_2,W_2)\), the row conjecture is
\((hW+tW)/2\), column's is \((HW+TW)/2\),
and third's is the \(3/8,1/8,1/8,3/8\) action-pair law.
To check mutual knowledge, row and column consider only
states in box 2 with third type \(W_1\) or \(W_2\);
both third types have the stated third conjecture.
The row and column types in box 2 all have the stated
row and column conjectures. Type \(W_2\) considers boxes
2 and 3; all row and column types in those boxes conjecture
fair opponent action and certainty of \(W\), and its own
conjecture is fixed by its type. Each player's entire
probability support therefore has the same full conjecture
profile. This is mutual knowledge of that profile.
The third conjecture's two marginals are fair, so individual
conjectures agree as stated.

For completeness, write row's probability of \(H\) as \(x\)
and column's probability of \(h\) as \(y\).
Row's payoffs to \(H,T\) are \(y,1-y\); column's payoffs
to \(h,t\) are \(1-x,x\), independent of the third action.
If \(y>1/2\), row must choose \(H\), so \(x=1\), which
makes column choose \(t\), contradicting \(y>1/2\).
If \(y<1/2\), the reverse forces \(x=0\) and column chooses
\(h\), again a contradiction. Thus \(y=1/2\).
If \(x>1/2\) column would choose \(t\), and if \(x<1/2\)
column would choose \(h\); hence \(x=1/2\).
Under their independent fair mixtures, matching occurs
with probability \(1/2\), so \(W\) yields \(3/2<2\).
The third player must choose \(E\). Conversely this profile
makes both first players indifferent and the third uniquely
optimal, proving existence and uniqueness.
In particular \((\text{fair},\text{fair},\delta_W)\) is not Nash.

Here is a full extension for an arbitrary positive integer \(m\).
Let \(L=2m+5\), and use boxes \(r=1,\ldots,L\).
Put \(w_r=3^{L-r}\), \(Z=4\sum_{r=1}^Lw_r\), and give
each cell of box \(r\) mass \(w_r/Z\).
Its diagonal cells have third type \(W_r\). The first
box's off-diagonal cells have type \(E_1\); the off-diagonal
cells of box \(r\geq2\) have type \(W_{r-1}\).
All other states receive mass zero.

Every row and column type again faces a fair opposing action.
For each \(r<L\), type \(W_r\) combines the diagonal of box
\(r\) and off-diagonal of box \(r+1\), whose weights have
ratio 3. It therefore has the same \(3/8,1/8,1/8,3/8\)
conjecture and strictly prefers \(W\). The endpoint types
\(E_1,W_L\) have the same rationality calculations as before.
All type marginals are positive, so all these are ordinary
conditional probabilities.

Let \(D\) be the event of the three fixed full conjectures
from the middle of the three-box example. Every occupied
state in boxes \(2,\ldots,L-1\) belongs to \(D\):
row and column conjecture certainty of \(W\), and every
third type encountered there is among \(W_1,\ldots,W_{L-1}\).
Starting at a positive-prior state, a belief edge remains
in the prior support. Row or column conditioning stays
in the same box; third-player conditioning moves by at
most one box. Start at the diagonal state
\((H_{m+3},h_{m+3},W_{m+3})\).
Every path of length at most \(m\) stays in boxes
\(3,\ldots,2m+3\), a subset of \(2,\ldots,L-1\).
By the finite path characterization, the state belongs
to \(K^rD\) for every \(1\leq r\leq m\).
It has positive prior mass, agreement of the stated marginals,
and the same non-Nash conjecture profile. This proves the
claimed finite-depth failure by explicit finite models.
\end{proof}
\end{annexitem}

\begin{annexitem}{annex:p17.ab.no.prior}{thm:p17-ab-no-prior}{All payoffs and incompatibility in Example 5.3}
\begin{proof}
Every row and column type faces the fair opponent mixture,
so both its actions yield \(1/2\). The sole third type
expects a matching pair with probability \(3/4\), giving
\(9/4\) from \(W\) and 2 from \(E\). All types are rational.
All types of a given player have identical action conjectures
and identical own payoffs. Thus the game, full conjecture
profile and rationality hold at every state and are commonly
known. The third conjecture has fair row and column marginals,
giving the agreed profile, whose non-Nash property was computed
in the preceding proof.

If a common prior existed, the third player's sole type
would have prior probability one. Its conditional identity
would force the joint prior on row and column types to have
the masses \(3/8,1/8,1/8,3/8\). Both row-type marginals
would equal \(1/2\). Conditional on row type \(H_1\),
column type \(h_1\) would then have probability \(3/4\),
contrary to that type's specified probability \(1/2\).
This contradicts (CP), so no compatible common prior exists.
\end{proof}
\end{annexitem}

\begin{annexitem}{annex:p17.ab.depth.disagreement}{thm:p17-ab-depth-disagreement}{A weighted finite path with different constant posteriors}
\begin{proof}
Let \(L=m+2\), and index \(2L+2\) occupied states by
\(k=0,\ldots,2L+1\). Give occupied state \(k\) probability
\(2^k/Z\), where \(Z=\sum_{k=0}^{2L+1}2^k\).
Player 1's types identify the pairs
\(\{0,1\},\{2,3\},\ldots,\{2L,2L+1\}\).
Player 2's types identify the singleton \(\{0\}\), the
pairs \(\{1,2\},\{3,4\},\ldots,\{2L-1,2L\}\), and the
singleton \(\{2L+1\}\). Player 3 has one type for each
individual \(k\). Represent occupied state \(k\) by its
three type coordinates; the third coordinate makes these
tuples distinct. Give all other tuples mass zero.

Players 1 and 2 have just one action each. Type \(k\) of
player 3 chooses 1 for odd \(k\) and 0 for even \(k\).
All payoffs are zero, and beliefs are conditional laws of
the specified prior. Every type has positive marginal
mass. Every player is rational, and the game and rationality
are commonly known.

In each player-1 pair the odd state's weight is twice the
even state's weight, so its conjecture assigns probability
\(2/3\) to player 3 choosing 1. In each nonsingleton
player-2 pair the even weight is twice the odd weight,
so its corresponding probability is \(1/3\).
Player 3's conjecture is constant because both opponents
have only one action. Hence the full conjecture profile
\(\varphi\) with these three laws is realized at every
occupied state \(1,\ldots,2L\).

A positive-probability belief transition stays in the
prior support and changes \(k\) by at most one; for player
3 it does not change \(k\). Start at \(k=m+2\).
All states reachable in at most \(m\) transitions have
indices between 2 and \(2m+2\), within \(1,\ldots,2L\).
The path characterization proves every \(K^r[\phi=\varphi]\)
through depth \(m\), while the two conjectures about player
3 remain \(2/3\ne1/3\). This proves the assertion with
normalized finite masses and no appeal to an infinite
agreement theorem.
\end{proof}
\end{annexitem}

\begin{annexitem}{annex:p17.ab.independent}{thm:p17-ab-independent}{Assumed independence and agreement}
\begin{proof}
At every state in the event asserted to be mutually known,
each full opponent law factors into its individual marginals
by the product lemma, applied to the opponents' finite
action coordinates. With one opponent this is automatic.
The asserted marginals make that law
\(\varphi^i=\bigotimes_{j\ne i}\sigma_j\).
Thus the original hypotheses imply mutual knowledge of this
fixed full-conjecture profile as well as of the game and
rationality. If \(\sigma_j(c_j)>0\), any other player's
conjecture has that positive marginal. The supported-action
lemma makes \(c_j\) optimal against
\(\bigotimes_{k\ne j}\sigma_k\). The Nash support criterion
finishes the proof. No prior or higher belief level enters.
\end{proof}
\end{annexitem}

\begin{annexitem}{annex:p17.ab.converse}{thm:p17-ab-converse}{Realization and the exact pointwise converse}
\begin{proof}
Given a Nash profile \(\sigma\), let \(S_i\) contain one
type for every \(c_i\) with \(\sigma_i(c_i)>0\). The type
chooses \(c_i\), has payoff function \(g_i^*\), and uses
the product of opponents' \(\sigma_j\) probabilities on
their corresponding types. Define \(P\) as the product
of all the \(\sigma_i\) type distributions.
Every type has positive marginal mass, and direct division
shows that the specified opponent laws are its conditional
distributions. Thus (CP) holds. Nash optimality makes
every type rational. Payoffs and full conjectures are
constant at all states, so all required mutual and common
knowledge events are the whole space. This proves the
existential realization and its positive-prior condition.
A pure equilibrium can use just one type per player.

For the pointwise assertion, suppose first that rationality
is mutually known. The assumed product conjectures and the
supported-action lemma give Nash of \(\sigma\), exactly as
in the proof of Remark 7.1.
Since \(R_i\) is an own-type event and \(K^1R\subseteq K_iR_i=R_i\),
each actual action is rational. Mutual knowledge of the
specified game and conjectures also makes their own-coordinate
components actual, by the same certainty reasoning.
Hence the actual actions are best responses to \(\sigma\).

Conversely suppose \(\sigma\) is Nash and every actual own
action is such a best response. Fix player \(i\).
On the probability-one event where the
specified game and full conjectures hold, \(i\)'s own
action remains its actual \(a_i(s_i)\), and is therefore
optimal against the stated opponent law. Furthermore the
event that an opponent action \(a_j\), \(j\ne i\), lies
outside the support of \(\sigma_j\) has \(i\)-probability zero:
\(\sigma_j\) is the corresponding marginal of the actual
product conjecture of \(i\). There are finitely many opponents,
so their actions all belong to these supports with probability one.
Every such supported action is a best response by Nash.
On the intersection of the stated probability-one events all
players are therefore rational. Thus \(p_i(R;s_i)=1\).
Repeating for every \(i\) gives \(K^1R\).
Positive equilibrium support implies best response by the
same support criterion, proving the stated sufficient repair.
The following explicit countermodel proves why the extra
actual-action condition cannot simply be omitted.
This argument is valid for every finite \(n\geq2\).
The factorization half of Theorem B supplies its product premise
in the common-prior case whenever the common-conjecture event
has positive probability, without first assuming rationality.

For the preliminary converse, mutual knowledge of \(a=\alpha\)
and \(g=g^*\) gives actual own actions and payoff arrays equal
to their specified values. It also makes each actual conjecture
the point mass on \(\alpha^{-i}\). Thus actual best-response
conditions and Nash of the corresponding degenerate profile
coincide, proving equivalence with actual rationality of all
players. Mutual knowledge of actions alone does not imply
mutual knowledge of the induced conjectures, so the preceding
conjecture converse cannot be used to strengthen this conclusion.
The source illustration has mutually known actions and a constant
game, with a pure equilibrium and actual rationality but without
mutual rationality. It supplies the exact separating counterexample.
\end{proof}
\end{annexitem}

\begin{annexitem}{annex:p17.ab.support.failures}{thm:p17-ab-support-failures}{Every omitted premise isolated}
\begin{proof}
First give each of two players actions \(0,1\) and payoff
equal to the player's own action. Give each two types,
also denoted \(0,1\), whose actions are those labels.
Every type assigns certainty to the opponent's type 1.
All payoff functions and full conjectures are constant.
At the actual state \((0,0)\), both players choose a strictly
inferior action and are irrational, but their conjectures
are \((\delta_1,\delta_1)\), the unique Nash equilibrium.
The game and conjectures are commonly known. The prior
concentrated at \((1,1)\) satisfies (CP): the type-1
conditionals agree and the type-0 kernels lie on marginal
null types. This is a model allowed by the source, which
does not demand positive mass of every actual state.
It refutes note~31's unqualified pointwise converse, not
the valid existential realization or either forward theorem.

For actual rather than mutual rationality, retain those
payoffs and action-types, but give every type a fair belief
on the opponent types and let \(P\) be their independent
uniform prior. At \((1,1)\), a state of probability \(1/4\),
both players are rational. The game and fair conjectures
are commonly known. Their fair mixed profile is not Nash,
since each player's action 0 is strictly inferior.
Each player assigns probability \(1/2\) to the opponent's
irrational type, so mutual rationality fails exactly as required.

For knowledge of own payoff without mutual knowledge of
the game, use the same uniform prior and fair beliefs.
Type \(r\in\{0,1\}\) of player \(i\) has utility
1 when the own action equals \(r\), and 0 otherwise,
independently of opponents; its actual action is \(r\).
Every type is rational, so rationality is commonly known.
Conjectures are still commonly known and fair, and each
type knows its own payoff function. At state \((0,0)\)
the actual game makes action 0 strictly dominant for both
players, so the fair conjecture profile is not Nash.
Each player considers an opponent payoff array favoring 1
possible, so the actual game is not mutually known.

Finally start with the complete model of Example 5.3.
For each player append a fresh anchor type. Keep every old
type's action, game and belief on old opponents unchanged.
Give the anchor types the same game, arbitrary actions,
and certainty of the opponents' anchor types. Put all prior
mass on the single all-anchor state. This is a common prior:
its positive-marginal anchor conditionals match the specified
kernels, and all old types have marginal probability zero.
On the product of the old types, the source game's rationality
and fixed full conjectures are commonly known, because every
old belief remains inside that product and all its types
are rational. Hence the event in (AB) for those conjectures
is nonempty, yet its conjecture profile is the non-Nash
\((\text{fair},\text{fair},\delta_W)\).
The all-anchor state has degenerate row/column conjectures,
so it is outside the specified full-conjecture event and
outside its common-knowledge event. The event in (AB)
therefore has prior probability zero. This explicitly
isolates the missing positivity rather than declaring
the given null-type kernels undefined.
\end{proof}
\end{annexitem}

\begin{annexitem}{annex:p17.ab.knowledge.of.equilibrium}{thm:p17-ab-knowledge-of-equilibrium}{The specified equilibrium event and the limits of iteration}
\begin{proof}
For a fixed mixed profile \(\sigma\), let \(L_\sigma\)
be the event that it is Nash in the current game \(g(s)\).
This is measurable: it is a finite conjunction of
measurable expected-payoff inequalities. If \(\sigma\)
is Nash in the specified \(g^*\), then
\([g=g^*]\subseteq L_\sigma\).
Mutual knowledge of that game, by monotonicity, gives
mutual knowledge of \(L_\sigma\). Theorems A and B supply
the premise that \(\sigma\) is Nash in \(g^*\).
The same reasoning applies to the preliminary observation
when mutual knowledge of the actual payoff array is added.

To show that the preliminary observation alone need not
give mutual knowledge of this fact, use the two-player
types \(0,1\), independent uniform prior and fair type
beliefs. Let every type choose action 0. At type 0,
payoff is 1 for own action 0 and 0 for own action 1.
At type 1 reverse these two payoffs. At state \((0,0)\)
both players are rational, and the action profile \((0,0)\)
is commonly known because actions are constant.
Each player nevertheless considers the opponent's type 1
possible. At such a state the fixed profile \((0,0)\)
is not Nash. Thus its Nash property is not mutually known.

For failure of higher-order knowledge even under Theorems
A and B, use the six-state prior of the source illustration,
with the same row labels \(C_1,D_1,D_2\) and column labels
\(c_1,d_1,d_2\), but give both players the action set
\(\{0,1\}\) and make every type choose 0.
Let the first two types of each player have the payoff
array that rewards own action 0 by 1 and own action 1 by 0.
Let the third type of each player have the reversed array.
All type beliefs are still the two-cell conditionals of
the six-state prior. All conjectures are the constant
point masses at opponent action 0.

At \((C_1,c_1)\), row considers column types \(c_1,d_1\);
column considers row types \(C_1,D_1\). All these types
have the first payoff array and choose its unique optimum.
Thus the specified game, rationality and the full conjectures
are mutually known. The fixed profile \((0,0)\) is Nash
in that game. Full conjectures are commonly known everywhere,
and the displayed state has positive prior probability,
so Theorem B's positive event exists as well.

However row considers \((C_1,d_1)\) possible, and there
column considers row type \(D_2\) possible. At
\((D_2,d_1)\), action 0 is strictly inferior for row,
so the fixed profile is not Nash. Therefore at
\((C_1,d_1)\) the Nash event is not mutually known.
Row gives this state probability \(1/2\) from the original
state, so second-order mutual knowledge of that event
fails there. This proves the exact higher-order limitation.
\end{proof}
\end{annexitem}
